
\textbf{Основные результаты}

\begin{enumerate}
    \item Проведён теоретический анализ трёх алгоритмов сортировки. Выведены формулы сложности (\ref{eq:bubble_comp}, \ref{eq:quick}, \ref{eq:quick_worst}, \ref{eq:merge}) и составлена сводная таблица~\ref{tab:complexity}, что соответствует данным из~\cite{knuth1998}.
    \item Выполнена корректная реализация каждого алгоритма на двух языках программирования — Python и C++.
    \item Разработана воспроизводимая методика тестирования с генерацией случайных массивов и многократными замерами.
    \item Проведены экспериментальные замеры для $n = 1000, 5000, 10000, 50000$. Результаты сведены в таблицы~\ref{tab:results} и~\ref{tab:bubble}.
    \item Визуализированы зависимости времени от размера массива (рисунки~\ref{fig:time_chart} и~\ref{fig:comparison}).
    \item Проанализированы различия в производительности: C++ быстрее Python в среднем в 6–7 раз для эффективных алгоритмов.
\end{enumerate}

\textbf{Практическая значимость}

Разработанные программные модули могут быть использованы в учебных целях для демонстрации работы алгоритмов сортировки. Полученные количественные данные помогают при выборе языка реализации для задач, требующих интенсивных вычислений.

\textbf{Направления дальнейших исследований}

В будущем возможно расширение работы в следующих направлениях:
\begin{itemize}
    \item добавление других алгоритмов (пирамидальная сортировка, сортировка Timsort);
    \item тестирование на разных типах данных (строки, структуры);
    \item сравнение с реализациями из стандартных библиотек (\texttt{std::sort} в C++ и \texttt{list.sort()} в Python).
\end{itemize}

Таким образом, цель работы достигнута, все задачи выполнены, полученные результаты являются достоверными и могут служить основой для дальнейших исследований в области сравнительного анализа производительности языков программирования.